\section{Microarchitecture}
\label{sec:microarchitecture}
The functional VHDL hierarchy contains 48 entities and five packages. A single
$8\times8$ tensor engine serves embedding, projections, feed-forward layers,
scoring, and classification. Registered readers provide row-vector operands;
tiled online attention maintains running state without a full score matrix.
The selector and compactor implement stable grouping, chronological retention,
and summaries. The sequencer selects four-block or block-two-exit execution.

The peripheral connects AXI-Lite control, framed 128-bit AXI streams, typed model
loading, bank arming and identity checks, input storage, output serialization,
and completed-inference counter snapshots. The logical model map is compacted
into three resident allocations of 90,112 64-bit words each. This allocation
is 2,162,688 bytes before other state and physical mapping overhead; it is not
a placed BRAM/URAM utilization result.

Core activation and QKV reads are registered. RMS variance, pooling, RTS
summaries, and final attention normalization use iterative division. Scalar
parameter and operand transactions currently serialize, so a functional
$8\times8$ array does not establish high sustained utilization.
\writingtask{RTL leads}{Add the actual hierarchy/dataflow figure, bank and port
schedule, buffer lifetimes, arithmetic widths and rounding points, pipeline
stages, reset/error behavior, and a measured breakdown of controller/RTS costs.
Separate implemented logic from vendor-generated PS/DMA integration.}
